% Options for packages loaded elsewhere
\PassOptionsToPackage{unicode}{hyperref}
\PassOptionsToPackage{hyphens}{url}
\documentclass[
  ignorenonframetext,
]{beamer}
\newif\ifbibliography
\usepackage{pgfpages}
\setbeamertemplate{caption}[numbered]
\setbeamertemplate{caption label separator}{: }
\setbeamercolor{caption name}{fg=normal text.fg}
\beamertemplatenavigationsymbolsempty
% remove section numbering
\setbeamertemplate{part page}{
  \centering
  \begin{beamercolorbox}[sep=16pt,center]{part title}
    \usebeamerfont{part title}\insertpart\par
  \end{beamercolorbox}
}
\setbeamertemplate{section page}{
  \centering
  \begin{beamercolorbox}[sep=12pt,center]{section title}
    \usebeamerfont{section title}\insertsection\par
  \end{beamercolorbox}
}
\setbeamertemplate{subsection page}{
  \centering
  \begin{beamercolorbox}[sep=8pt,center]{subsection title}
    \usebeamerfont{subsection title}\insertsubsection\par
  \end{beamercolorbox}
}
% Prevent slide breaks in the middle of a paragraph
\widowpenalties 1 10000
\raggedbottom
\AtBeginPart{
  \frame{\partpage}
}
\AtBeginSection{
  \ifbibliography
  \else
    \frame{\sectionpage}
  \fi
}
\AtBeginSubsection{
  \frame{\subsectionpage}
}
\usepackage{iftex}
\ifPDFTeX
  \usepackage[T1]{fontenc}
  \usepackage[utf8]{inputenc}
  \usepackage{textcomp} % provide euro and other symbols
\else % if luatex or xetex
  \usepackage{unicode-math} % this also loads fontspec
  \defaultfontfeatures{Scale=MatchLowercase}
  \defaultfontfeatures[\rmfamily]{Ligatures=TeX,Scale=1}
\fi
\usepackage{lmodern}
\ifPDFTeX\else
  % xetex/luatex font selection
\fi
% Use upquote if available, for straight quotes in verbatim environments
\IfFileExists{upquote.sty}{\usepackage{upquote}}{}
\IfFileExists{microtype.sty}{% use microtype if available
  \usepackage[]{microtype}
  \UseMicrotypeSet[protrusion]{basicmath} % disable protrusion for tt fonts
}{}
\makeatletter
\@ifundefined{KOMAClassName}{% if non-KOMA class
  \IfFileExists{parskip.sty}{%
    \usepackage{parskip}
  }{% else
    \setlength{\parindent}{0pt}
    \setlength{\parskip}{6pt plus 2pt minus 1pt}}
}{% if KOMA class
  \KOMAoptions{parskip=half}}
\makeatother
\usepackage{longtable,booktabs,array}
\usepackage{caption}
\captionsetup[table]{skip=6pt}
\newcounter{none} % for unnumbered tables
\usepackage{calc} % for calculating minipage widths
\usepackage{caption}
% Make caption package work with longtable
\makeatletter
\def\fnum@table{\tablename~\thetable}
\makeatother
\setlength{\emergencystretch}{3em} % prevent overfull lines
\providecommand{\tightlist}{%
  \setlength{\itemsep}{0pt}\setlength{\parskip}{0pt}}
% Auto-generated by infrastructure/rendering/_pdf_latex_helpers.py::extract_math_font_preamble.
% Minimal header passed to Pandoc via -H so Beamer slides pick up the same math font
% as the combined-PDF path. Adjust the manuscript's preamble.md to change the font globally.
\usepackage{unicode-math}
\setmathfont{latinmodern-math.otf}

% Slide layout defaults for warning-clean scientific decks.
\usepackage{etoolbox}
\IfFileExists{xurl.sty}{\usepackage{xurl}}{}
\IfFileExists{seqsplit.sty}{\usepackage{seqsplit}}{\newcommand{\seqsplit}[1]{#1}}
\protected\def\breakseq#1{\seqsplit{#1}}
\protected\def\breaktt#1{\begingroup\ttfamily\seqsplit{#1}\endgroup}
\setlength{\emergencystretch}{6em}
\tolerance=5000
\hbadness=10000
\hfuzz=1pt
\setlength{\tabcolsep}{2pt}
\AtBeginEnvironment{longtable}{\tiny\renewcommand{\arraystretch}{0.86}\setlength{\tabcolsep}{1pt}}
\AtBeginEnvironment{tabular}{\tiny\renewcommand{\arraystretch}{0.86}\setlength{\tabcolsep}{1pt}}
\AtBeginEnvironment{equation}{\tiny}
\AtBeginEnvironment{equation*}{\tiny}
\AtBeginEnvironment{align}{\tiny}
\AtBeginEnvironment{align*}{\tiny}
\AtBeginEnvironment{itemize}{\footnotesize}
\AtBeginEnvironment{enumerate}{\footnotesize}
\AtBeginEnvironment{description}{\footnotesize}
\setbeamerfont{caption}{size=\tiny}
\setbeamerfont{caption name}{size=\tiny}
\setbeamerfont{normal text}{size=\small}
\setbeamerfont{frametitle}{size=\small}
\setbeamerfont{section title}{size=\footnotesize}
\setbeamerfont{subsection title}{size=\footnotesize}
\setbeamertemplate{section page}{%
  \centering
  \begin{beamercolorbox}[sep=12pt,center,wd=\paperwidth]{section title}
    \parbox{0.86\paperwidth}{\centering\usebeamerfont{section title}\insertsection\par}
  \end{beamercolorbox}
}
\setbeamertemplate{subsection page}{%
  \centering
  \begin{beamercolorbox}[sep=8pt,center,wd=\paperwidth]{subsection title}
    \parbox{0.86\paperwidth}{\centering\usebeamerfont{subsection title}\insertsubsection\par}
  \end{beamercolorbox}
}
\setlength{\abovecaptionskip}{2pt}
\setlength{\belowcaptionskip}{0pt}

% Natbib and cross-reference fallbacks — slides don't load natbib
% or cleveref, but combined-PDF manuscript prose may emit these
% commands. The fallback renders citations as a bracketed key list
% and cross-references as detokenized labels so slides don't fail on
% undefined control sequences or raw underscores. \providecommand is
% a no-op if packages load later.
\providecommand{\citep}[1]{[#1]}
\providecommand{\citet}[1]{#1}
\providecommand{\citealp}[1]{#1}
\providecommand{\citeauthor}[1]{#1}
\providecommand{\citeyear}[1]{#1}
\providecommand{\cref}[1]{\texttt{\detokenize{#1}}}
\providecommand{\Cref}[1]{\texttt{\detokenize{#1}}}

\newtheorem{proposition}{Proposition}
\newtheorem{hypothesis}{Hypothesis}
\usepackage{bookmark}
\IfFileExists{xurl.sty}{\usepackage{xurl}}{} % add URL line breaks if available
\urlstyle{same}
% fallback for those not using the hyperref driver hyperxmp:
\makeatletter
\@ifundefined{xmpquote}{\newcommand{\xmpquote}[1]{#1}}{}
\makeatother
\hypersetup{
  hidelinks,
  pdfcreator={LaTeX via pandoc}}

\author{\texorpdfstring{}{}}
\date{}

\begin{document}

\section{Related Work: Historical Foundations, Active Inference Evidence
Status, and Information
Sources}\label{sec:related-work-aif-information-sources}

\subsection{Active Inference and FEP lineage: historical context and
curriculum
scope}\label{active-inference-and-fep-lineage-historical-context-and-curriculum-scope}

\begin{frame}[allowframebreaks]{Formal lineage and model commitments}
\protect\phantomsection\label{formal-lineage-and-model-commitments}
The FEP is commonly situated in a perception-as-inference lineage that
includes Helmholtz, Bayesian-brain accounts of uncertainty, and tutorial
treatments that translate the mathematics into worked learning sequences
(Helmholtz 1867; Knill and Pouget 2004; Bogacz 2017; Sprevak and Smith
2023). Its modern variational formulation is developed in Friston's
early and review papers (Friston et al. 2006; Friston 2010), and recent
overview work usefully maps the breadth of related research without
removing the need for local modeling commitments (Friston et al. 2023;
Zhang and Xu 2024). Mathematical reviews and textbooks situate the
principle as a variational bound on surprise and a process theory for
perception, action, and learning (Buckley et al. 2017; Parr et al. 2022;
Smith et al. 2022; Pezzulo et al. 2024). Associative-learning work makes
\end{frame}

\begin{frame}[allowframebreaks]{Formal lineage and model commitments}
attention, prediction-error, blocking, and surprise phenomena explicit
inside an Active Inference account, which is useful pedagogical context
but not direct evidence for this curriculum's outcomes (Anokhin et al.
2024). The discrete POMDP formulation with \(A\), \(B\), \(C\), and
\(D\) matrices is the direct basis for the kernels here (K. Friston et
al. 2017; Da Costa et al. 2020), while epistemic value and curiosity are
treated by the expected-free-energy literature (Friston et al. 2015; K.
J. Friston, Lin, et al. 2017; Millidge, Tschantz, et al. 2021; Sajid et
al. 2021; Champion et al. 2026). Recent variational-inference and
message-passing treatments make those formulation choices explicit by
separating epistemic-prior, entropy-correction, and planning-correction
assumptions rather than treating EFE as one unqualified objective
(Nuijten, Lukashchuk, et al. 2026, 2025; Nuijten, Laar, et al. 2026).
\end{frame}

\begin{frame}[allowframebreaks]{Formal lineage and model commitments}
The planning-frontier argument therefore leans on those EFE planning and
message-passing papers rather than on informal Lazy Dynamics talk
material, because the manuscript needs source-backed claims about
objectives, factorization, and assumptions rather than a weak
conference-talk framing.
\end{frame}

\subsection{Ancient, medieval, and early-modern foundations of
inference, perception, and sequenced
learning}\label{ancient-medieval-and-early-modern-foundations-of-inference-perception-and-sequenced-learning}

\begin{frame}[allowframebreaks]{Ancient, medieval, and early-modern
foundations of inference, perception, and sequenced learning}
This is a selective history of problems and design concerns, not a
single-lineage account of Active Inference. The sources below are used
to motivate an inspectable prerequisite order and to distinguish
historical analogies from modern formal commitments. They do not show
that modern FEP, Bayesian inference, variational inference, adaptive
tutoring, or Active Inference existed before their later mathematical
and computational formulations. The historical source set spans
demonstration, optics, method, probability, association, education, and
structured cognition (Aristotle 1901; Ibn al-Haytham 1989; Quintilian 95
AD; Bacon 1620; Descartes 1628; Huygens 1657; Bernoulli 1713; Moivre
1718; Bayes and Price 1763; Laplace 1774; Newton 1704; Berkeley 1709;
Locke 1690, 1693; Hume 1739, 1748; Hartley 1749; Condillac 1754;
Comenius 1657; Rousseau 1762; Kant 1781).
\end{frame}

\begin{frame}[allowframebreaks]{Demonstration, vision, and ordered
instruction before modern probability}
\protect\phantomsection\label{demonstration-vision-and-ordered-instruction-before-modern-probability}
Ancient demonstration and explanation provide one early vocabulary for
separating reasons, evidence, and teachable order (Aristotle 1901).
Quintilian's \emph{Institutio Oratoria} presents a staged rhetorical
education involving elementary studies, practice, revision, and
progression through increasingly demanding forms of performance
(Quintilian 95 AD). Its relevance here is historical: it shows that
ordered instruction and repeated practice were explicit design concerns,
not that ancient rhetoric implemented a learner model or adaptive tutor.

\end{frame}

\begin{frame}[allowframebreaks]{Demonstration, vision, and ordered
instruction before modern probability}
The medieval Arabic optical tradition adds a distinct perception strand.
Ibn al-Haytham's \emph{Book of Optics}, represented here through a
critical edition and translation, treats vision through geometrical
analysis, observation, and experimentally constrained accounts of seeing
(Ibn al-Haytham 1989). That tradition belongs in the history of
perception and evidence, but it is not a predictive-coding theory, a
Bayesian brain model, or evidence for the biological adequacy of the
kernels in this release.
\end{frame}

\begin{frame}[allowframebreaks]{Early-modern method, probability, and
experimental perception}
\protect\phantomsection\label{early-modern-method-probability-and-experimental-perception}
Early-modern method sources connect ordered reasoning to observation
without collapsing their differences. Bacon's induction and experimental
program, and Descartes's rules for moving through clear intuitions,
deductions, and probable conjectures, make method and sequence explicit
(Bacon 1620; Descartes 1628). Huygens's work on expectation in games of
chance provides an early published probability treatment, while
Bernoulli, de Moivre, Bayes, and Laplace extend the mathematical history
of repeated observation and inverse-probability reasoning (Huygens 1657;
Bernoulli 1713; Moivre 1718; Bayes and Price 1763; Laplace 1774). These
sources support a history of calculation and inference under
uncertainty, not a claim that their concepts were already variational or
computational in the modern sense.

\end{frame}

\begin{frame}[allowframebreaks]{Early-modern method, probability, and
experimental perception}
Newton's \emph{Opticks} adds a controlled-experiment and
sensory-observation strand: optical phenomena are varied, separated, and
analyzed through explicit apparatus and mathematical argument (Newton
1704). Berkeley's account of vision similarly treats visual perception
as involving learned relations among sensory cues, but neither source
should be translated directly into modern generative models or neural
predictive coding (Berkeley 1709).
\end{frame}

\begin{frame}[allowframebreaks]{Eighteenth-century association,
education, and structured cognition}
\protect\phantomsection\label{eighteenth-century-association-education-and-structured-cognition}
Eighteenth-century sources then widen the educational and psychological
context. Locke, Hume, Hartley, and Condillac discuss experience,
association, habit, sensation, and reflection in different philosophical
registers (Locke 1690; Hume 1739, 1748; Hartley 1749; Condillac 1754).
Comenius, Locke's education writings, and Rousseau's \emph{Émile} make
staged progression, practice, motivation, and learner-sensitive ordering
visible as historical pedagogical concerns (Comenius 1657; Locke 1693;
Rousseau 1762). Kant provides a different account of structured
cognition and recognition (Kant 1781). Together these sources justify
treating sequence, representation, and evidence status as design
questions; they do not establish modern cognitive science, personalized
learning, or empirical efficacy.
\end{frame}

\subsection{Evidence-status critiques and empirical
boundaries}\label{evidence-status-critiques-and-empirical-boundaries}

\begin{frame}[allowframebreaks]{Evidence-status critiques and empirical
boundaries}
The same literature also warns against collapsing every layer of the
project into a single grand claim. A variational identity, a process
theory, a simulation recipe, and an empirical neuroscience hypothesis
have different evidential burdens. Technical critiques of the FEP's
particular-physics assumptions and subsequent Bayesian-mechanics framing
therefore belong in the curriculum as boundary-setting material, not as
adversarial footnotes (Aguilera et al. 2022, 2023). Recent empirical
reviews similarly motivate the platform's wording discipline: the
curriculum can teach predictive coding and Active Inference as formal
frameworks while still marking where direct empirical support is
incomplete or domain-specific (Hodson et al. 2024).
Psychology-and-psychiatry reviews add the same caution at the
behavioral-modeling level: Active Inference is productive for model
fitting and explanation, but its theory-level empirical adequacy still
\end{frame}

\begin{frame}[allowframebreaks]{Evidence-status critiques and empirical
boundaries}
requires direct comparison with alternatives (Badcock and Davey 2024).
Gershman's critique is useful here because it asks what distinctive
predictions the FEP contributes, which is exactly the kind of scope
boundary this manuscript treats as a claim-support obligation (Gershman
2019). The perception literature supplies a complementary, narrower
evidence map: Bayesian vision and predictive-processing accounts connect
generative hypotheses to perceptual inference, while
cortical-microcircuit, efficient-coding, perceptual-decision, and
neurophysiological reviews expose model-specific assumptions and
unresolved empirical questions (Yuille and Kersten 2006; Clark 2013;
Hohwy 2013; Kersten et al. 2004; Bastos et al. 2012; Spratling 2017;
Summerfield and Lange 2014; Wei and Stocker 2015; Millidge, Seth, et al.
2021; Smith et al. 2020; Seth and Hohwy 2020). These sources sharpen the
perception-as-inference lane without turning predictive processing into
\end{frame}

\begin{frame}[allowframebreaks]{Evidence-status critiques and empirical
boundaries}
a settled theory of the brain or this curriculum into a neuroscience
intervention.
\end{frame}

\subsection{Information, variational, filtering, and application
sources}\label{information-variational-filtering-and-application-sources}

\begin{frame}[allowframebreaks]{Information and filtering sources}
\protect\phantomsection\label{information-and-filtering-sources}
The information-theoretic and variational strata follow standard
treatments of entropy, KL divergence, variational inference,
graphical-model message passing, the ELBO/free energy relationship, and
rate-distortion or bottleneck views of task-relevant compression
(Shannon 1948; Beal 2003; Cover and Thomas 2006; Tishby et al. 1999;
Bishop 2006; Jordan et al. 1999; Winn and Bishop 2005; Blei et al.
2017). State-space filtering, continuous-time predictive coding, and
neural message passing are represented in the curriculum, with deeper
executable worked examples left for later releases; the relevant
literature includes Bayesian filtering, Helmholtz machines, hierarchical
predictive coding, the graphical-brain account, DEM, generalized
filtering, and generalized free energy (Särkkä 2013; Dayan et al. 1995;
Rao and Ballard 1999; Friston 2008; K. J. Friston, Parr, et al. 2017;
\end{frame}

\begin{frame}[allowframebreaks]{Information and filtering sources}
Friston et al. 2008; Friston et al. 2010; Parr and Friston 2019; Kataoka
and Doya 2026; Tschantz et al. 2023). Probability, information, and
Bayesian-analysis textbooks provide adjacent mathematical foundations
for these implementations and help distinguish a reusable inference
method from any one Active-Inference formulation (MacKay 2003; Jaynes
2003; Murphy 2012, 2022; Gelman et al. 2013).

Broader applications to embodied cognition, biological
self-organization, collective behavior, action-oriented model learning,
and institutional sensemaking motivate the applications and practice
strata (Clark 2016; Friston 2013, 2019; Tschantz et al. 2020; Ramstead
et al. 2018). Application-facing work spans cognitive control,
hierarchical motivation, computational consciousness, clinical
neuroscience, control-as-inference comparisons, robotics, and
model-based cognitive flexibility (Pezzulo 2012; Pezzulo et al. 2018,
2023; Vilas et al. 2022; Smith et al. 2021; Imohiosen et al. 2020;
\end{frame}

\begin{frame}[allowframebreaks]{Information and filtering sources}
Millidge et al. 2020; Schwartenbeck et al. 2015; Lanillos et al. 2021;
Catal et al. 2021; Queißer et al. 2021; Sales et al. 2019). These are
application and comparison surfaces, not evidence that the present
discrete kernels scale to robotics, clinical decision support, or human
cognitive change. Predictive-learning neuroimaging evidence strengthens
the plausibility of predictive learning as a representational mechanism,
but it remains neural-task evidence rather than evidence that this
microlearning curriculum changes learner representations (Greco et al.
2024). Current frontier sources sharpen the roadmap boundary. Reactive
message passing and RxInfer show a software and algorithmic path for
scalable Bayesian message passing; coupled free energy points toward
robust variational geometry; active predictive coding and
divide-and-conquer predictive coding extend the neural-inference and
scalability context. Those sources are cited as scope-extending
\end{frame}

\begin{frame}[allowframebreaks]{Information and filtering sources}
scholarship, not as implemented kernels, production-scale evidence, or
learner-outcome evidence in this release (Bagaev and Vries 2023; Bagaev
et al. 2023; Nelson et al. 2025; Rao et al. 2024; Sennesh et al. 2024).
Recent experimental work also grounds the FEP outside purely theoretical
exposition (Isomura et al. 2023).
\end{frame}

\end{document}
